\section{Limitations}

\paragraph{Hybrid human-machine generated data.} Our dataset is sourced primarily from text generated by LLMs. We pursued this method, which has quickly emerged as a popular alternative to time-intensive manual data collection \cite{kim-etal-2023-soda, jang2023conversation}, to avoid the logistical and legal complexities of collecting very long-term real-world conversations at scale. We ensure that the dataset mirrors real-world interactions as much as possible by having human annotators verify and edit the generated conversations. However, we acknowledge that this dataset may not fully reflect the nuances of real-world online conversations. 

\paragraph{Limited exploration of multimodal behavior.} Since the images in our dataset are sourced from the web, they do not demonstrate the visual long-term consistencies that are usually exhibited in personal photos (e.g., appearance, home environment, people and pets, etc.). Consequently, we find that the images in our dataset can be replaced with their captions without much loss of information, except for cases where OCR is required. Nevertheless, our work is a first step toward research into the multimodal aspect of very long-term conversations.

\paragraph{Language.} Our LLM-based pipeline for generating long-term conversations has been developed for the English language only. However, our pipeline can be made to work with any other language using an LLM that is proficient at that language and appropriate translations of our prompts.

\paragraph{Closed-source LLMs.} We use state-of-the-art LLMs in our dialog generation pipeline to create a dialog dataset that is as realistic as possible. Unfortunately, this meant employing the strongest commercial LLMs available through a paid API, similar to many concurrent works that generate synthetic conversations \cite{zhong2023memorybank, lu2023memochat}. We will make the code for our generative pipeline publicly available in the hope that it can be made to work effectively with state-of-the-art open-source LLMs in the future.

\paragraph{Evaluation of long-form NLG.} LLMs are prone to generating verbose answers even when prompted to answer in short phrases. This creates challenges in evaluating the correctness of answers provided by LLMs and has been widely documented in NLP literature \cite{chang2023booookscore, xu2023critical, krishna2023longeval}. Our evaluation framework suffers from the same challenges when used for experimenting with LLMs.
